\thesisappendix{Derived geometric and elastic quantities}{app:derived-quantities}

Every derived quantity used in the v2 contract --- temperature
perturbation $\theta$, propagation vector $d = x_p - x_s$, distance,
unit direction, azimuth, displacement magnitude, Lam\'e parameters
--- is computed by a single Python module that acts as the
authoritative source of truth for these formulas in the prototype.
The module is pure Python (no FastAPI or Pydantic dependencies), and
each formula is pinned to a reference value by a unit test in the
backend test suite.

\begin{lstlisting}[language=Python,caption={Derived-quantities module (excerpt)}]
REFERENCE_TEMPERATURE_K = 273.15
SOURCE_TEMPERATURE_K = 1500.0
DOMAIN_SIZE_M = 1.0


def compute_theta_k(
    source_temperature_k: float = SOURCE_TEMPERATURE_K,
    reference_temperature_k: float = REFERENCE_TEMPERATURE_K,
) -> float:
    """Temperature perturbation theta = T_source - T_ref (kelvin)."""
    return float(source_temperature_k) - float(reference_temperature_k)


def compute_distance_m(source, probe):
    dx = probe[0] - source[0]
    dy = probe[1] - source[1]
    return math.hypot(dx, dy)


def compute_azimuth_deg(source, probe):
    """Propagation azimuth in degrees, atan2 convention, xy-plane."""
    dx, dy = probe[0] - source[0], probe[1] - source[1]
    return math.degrees(math.atan2(dy, dx))


def compute_displacement_magnitude_m(u_m, v_m):
    return math.hypot(u_m, v_m)


def derive_elastic_constants(
    young_modulus_pa, poisson_ratio,
    rho_kg_m3=None, vp_m_s=None, vs_m_s=None,
    heat_capacity_j_kgk=None,
):
    """Prefers (E, nu); falls back to (rho, Vp, Vs) under
    isotropic-elasticity assumptions."""
    if young_modulus_pa is not None and poisson_ratio is not None:
        E, nu = float(young_modulus_pa), float(poisson_ratio)
        if nu == 0.5:
            raise ValueError("Poisson ratio nu = 0.5 is degenerate.")
        shear = E / (2.0 * (1.0 + nu))
        bulk = E / (3.0 * (1.0 - 2.0 * nu))
        lame_lambda = E * nu / ((1.0 + nu) * (1.0 - 2.0 * nu))
    elif rho_kg_m3 is not None and vp_m_s is not None and vs_m_s is not None:
        rho, vp, vs = float(rho_kg_m3), float(vp_m_s), float(vs_m_s)
        shear = rho * vs * vs
        bulk = rho * (vp * vp - (4.0 / 3.0) * vs * vs)
        lame_lambda = bulk - (2.0 / 3.0) * shear
    else:
        raise ValueError("provide (E, nu) or (rho, Vp, Vs).")
    return DerivedElasticConstants(
        shear_modulus_pa=shear,
        bulk_modulus_pa=bulk,
        lame_lambda_pa=lame_lambda,
        volumetric_heat_capacity_j_m3k=(
            float(rho_kg_m3) * float(heat_capacity_j_kgk)
            if rho_kg_m3 is not None and heat_capacity_j_kgk is not None
            else None
        ),
    )


def verify_catalog_elastic_consistency(
    young_modulus_pa, derived_shear_pa, derived_bulk_pa,
    *, relative_tolerance=0.01,
):
    """Sanity check: E should match 9 K mu / (3 K + mu) within 1%."""
    denom = 3.0 * derived_bulk_pa + derived_shear_pa
    if denom <= 0:
        return False
    e_implied = 9.0 * derived_bulk_pa * derived_shear_pa / denom
    return abs(e_implied - young_modulus_pa) / abs(young_modulus_pa) \
        < relative_tolerance
\end{lstlisting}

The closed-form fallback paths in \texttt{derive\_elastic\_constants}
match the analytical formulas of Section~2.6: $\mu = \rho V_s^{2}$,
$K = \rho(V_p^{2} - \tfrac{4}{3} V_s^{2})$, and $\lambda = K -
\tfrac{2}{3}\mu$. The startup-time check
\texttt{verify\_catalog\_elastic\_consistency} guards the v2 catalog
against drift between the stored $E$ and the
$(\rho, V_p, V_s)$ triple that was used to derive it.
